\documentclass[10pt,a4paper]{amsart}
\usepackage[margin=19mm]{geometry}
\usepackage{amsmath,amssymb,mathtools}
\usepackage[scr=rsfs,cal=euler]{mathalfa}
\usepackage{xfrac}
\usepackage[unicode,hidelinks,bookmarks=false]{hyperref}
\newcommand{\ie}{i.e.\ }
\newcommand{\cf}{cf.\ }
\newcommand{\resp}{respectively\ }
\newcommand{\Subsection}{\subsection}
\providecommand{\nobreakdash}{\nobreak\hbox{-}}
\setlength{\emergencystretch}{2em}
\multlinegap0pt
\usepackage{amssymb}
\newtheorem{proposition}{Proposition}
\newtheorem{assumption}{Assumption}
\newtheorem{lemma}{Lemma}
\newtheorem{theorem}{Theorem}
\title{Resonance-based integrators for stochastic Schrödinger equations. Convergence and long-time error bounds.}
\author{Stefano Di Giovacchino}
\date{Source: arXiv:2410.22201v3 (CC BY 4.0)}
\begin{document}
\maketitle

\noindent\emph{Source and licence.} Resonance-based integrators for stochastic Schrödinger equations. Convergence and long-time error bounds.. Version arXiv:2410.22201v3 (CC BY 4.0), distributed by arXiv under the Creative Commons Attribution 4.0 International licence, http://creativecommons.org/licenses/by/4.0/, as recorded in the arXiv licence metadata for this exact version and shown on the abstract page https://arxiv.org/abs/2410.22201v3. Full licence notice: \texttt{licenses/arxiv-2410.22201-CC-BY-4.0.txt}.\par\smallskip

\noindent\textit{Editorial scope.} Counted unit: the article's theorem long\_term\_path\_orig\_cub (source0.tex lines 1981--1988). The article's complete original proof of it is delivered with it (source0.tex lines 1989--2064, one \texttt{proof} environment of 76 lines). No mathematics is written, repaired or re-derived: the statement and the proof are the article's own lines, in the article's own notation, and no retained line is substituted or re-flowed.  Retained uncounted as the source-local context the proof needs: lines 1340--1346; lines 1402--1405; lines 1409--1416; lines 1455--1458; lines 1463--1476.  Nothing was omitted from any retained span, so the package carries no omission record.  No retained line contains a citation, so the wrapper carries no bibliography and no bibliography entry.  No retained line contains a figure environment, an image-inclusion command or drawing code, so no image is delivered and the package declares no asset dependency.  Editorial formatting only, in addition: an AMS \texttt{amsart} preamble that restates the article's own declarations and macros used by the retained lines, namely the environments \texttt{proposition}, \texttt{assumption}, \texttt{lemma}, \texttt{theorem} on separate counters, together with the standard packages those lines need.  The retained environments print in this document as Proposition 1, Assumption 1, Lemma 1, Theorem 1 in order of appearance, whereas the article prints the same items with the numbering of its own full document.  The complete native source of the article as distributed in the arXiv e-print for this exact version is bundled unchanged as \texttt{sources/167603/source0.tex}.
    \begin{proposition}
    \label{conv_path1}
        Let us assume $\psi_0\in H^{\sigma+\gamma}$ and $Q^{\frac{1}{2}}\in\mathcal{L}_2^{\sigma+\gamma}$, for $\sigma\in \mathbb{N}^+$ and $\gamma\in(0,1]$. Then, for any $\delta<\gamma$, there exists a random variable $K_\delta(\omega)$ such that one has
        $$
        \max_{1\le n\le N}\|\psi(t_n)-\psi_n\|_{H^\sigma}\le K_\delta(\omega) \tau^\delta,\qquad \mathbb{P}-{\rm a.s.}
        $$
    \end{proposition}

\begin{equation}
    \label{cub_small}
    d\psi^\varepsilon(t)={\bf i}\Delta \psi^\varepsilon(t) dt-{\bf i}\varepsilon^2 |\psi^\varepsilon(t)|^2\psi^\varepsilon(t)dt+\varepsilon Q^{\frac{1}{2}}dW(t), \qquad \psi^\varepsilon(0)=\phi(x).
\end{equation}

\begin{assumption}
\label{smal_cub_ass}
For $\phi\in H^\sigma$ and $Q^{\frac{1}{2}}\in\mathcal{L}_2^\sigma$, for $\sigma\ge 2$, we assume the following bound for the exact solution to  \eqref{cub_small}
$$
\big(\mathbb{E}\big[\displaystyle\sup_{t\in[0,T_\varepsilon]}\|\psi^\varepsilon(t)\|^{2p}_{H^\sigma}\big]\big)^{\frac{1}{2p}}\lesssim 1, \qquad T_\varepsilon=\frac{T}{\varepsilon^2},
$$
for any $p\in\mathbb{N}^+$.
\end{assumption}

\begin{equation}
    \label{trunc_cub2}
 d\psi^{\varepsilon,R}={\bf i}\Delta \psi^{\varepsilon,R} dt-{\bf i}\varepsilon^2 \theta_R(\psi^{\varepsilon,R})|\psi^{\varepsilon,R}|^2\psi^{\varepsilon,R} dt+\varepsilon Q^{\frac{1}{2}}dW, \qquad \psi^{\varepsilon,R}(0)=\phi(x),
\end{equation}

\begin{lemma}
\label{long_term_cub_trunc}
    Let us assume Assumption \ref{smal_cub_ass}. For any fixed $\tau_0\in(0,1)$ independent on $\varepsilon$, there exists a constant $K(\tau_0)$ such that, for any $\tau\in(0,K(\tau_0))$, the numerical integrator (SNRLR2), applied to the truncated cubic equation \eqref{trunc_cub2}, satisfies the following error estimate
    \begin{equation}
        \label{long_term_mean_trunc_cub}
    \big(\mathbb{E}\big[\max_{1\le n\le N_\varepsilon}\|\psi^{\varepsilon,R}(t_n)-\psi^{\varepsilon,R}_n\|^2_{H^1}\big]\big)^{\frac{1}{2}}\lesssim\tau_0^{\sigma-1}+\varepsilon^2\tau, \qquad N_\varepsilon:=\frac{T}{\varepsilon^2 \tau}.
    \end{equation}
    In addition, if we assume a setting in which the term $\tau_0^{\sigma-1}$ in \eqref{long_term_mean_trunc_cub} is neglectable, there exists a sequence $\{\tau_k\}_{k\in\mathbb{N}^+}$, with $\tau_k\to 0, \ k\to \infty$, such that, for any $R>0$ and $\delta\in(0,1)$, one has
    \begin{equation}
        \label{long_term_path_trunc_cub}
        \max_{1\le n\le N^k_\varepsilon}\|\psi^{\varepsilon,R}(t_n)-\psi^{\varepsilon,R}_n\|_{H^1}<\varepsilon^2 \tau_k^\delta, \quad \forall k>k_0(\omega), \qquad \mathbb{P}-{\rm a.s.},
    \end{equation}
    where $N_\varepsilon^k=\frac{T}{\varepsilon^2 \tau_k}$, with $k_0(\omega)$ independent on $\varepsilon$.
\end{lemma}

\begin{theorem}
\label{long_term_path_orig_cub}
    Let us assume the setting of Lemma \ref{long_term_cub_trunc}, under the condition that the term $\tau_0^{\sigma-1}$ is neglectable. Then, the numerical integrator (SNRLR2), applied to \eqref{cub_small} satisfies 
    $$
    \max_{1\le n \le N_\varepsilon}\| \psi^\varepsilon(t_n)-\psi^\varepsilon_n\|_{H^1}<\varepsilon^2\tau^\delta, \qquad \mathbb{P}-{\rm a.s.}, \quad N_\varepsilon=\frac{T}{\varepsilon^2 \tau},
    $$
    for any $\delta<1$, asymptotically in $\tau$.
\end{theorem}

\begin{proof}
    Similarly to the proof of Proposition \ref{conv_path1}, let us denote
$$
R_{1,\varepsilon}(\omega)\ge \sup_{t\in[0,T_\varepsilon]} \|u(t)\|_{H^1}+\max_{1\le n\le N_\varepsilon} \|\mathcal{W}_n\|_{H^1},
$$
where we have explicit the fact that, in principle, one cannot assume that the random variable $R_{1,\varepsilon}$ is independent from $\varepsilon$.

     Now, let $\tau$ sufficiently small (independently on $\varepsilon$) and assume there exists $1\le \tilde{n}\le N_\varepsilon$ such that one has
     $$
         \|\psi^\varepsilon (t_n)-\psi^\varepsilon_{n}\|_{H^1}<\varepsilon^2 \tau^\delta, \ n=1,\dots,\tilde{n}-1, \quad \|\psi^\varepsilon (t_{\tilde{n}})-\psi^\varepsilon_{\tilde{n}}\|_{H^1}\ge \varepsilon^2 \tau^\delta.
     $$
     Then, for $1\le n\le \tilde{n}-1$ and considering $\varepsilon^2 \tau^\delta<1$, one has
     \begin{align*}
         \|\psi^\varepsilon_n\|_{H^1}&\le \|\psi^\varepsilon(t_n)\|_{H^1}+\|\psi^\varepsilon(t_n)-\psi^\varepsilon_n\|_{H^1}\\
         &\le \|\psi^\varepsilon(t_n)\|_{H^1}+\varepsilon^2 \tau^\delta\\
         &\le R_{1,\varepsilon}(w)+1.
     \end{align*}
     Then, using this, the triangular inequality, the unitary property of the operator ${\rm e}^{{\bf i}t\Delta}$ and the algebra property of $H^\sigma$, one has
     \begin{align*}
         \|\psi^\varepsilon_{\tilde{n}}\|_{H^1}\le R_{1,\varepsilon}(w)+1&+\varepsilon^2\tau C_{GN} \|\psi^\varepsilon_{\tilde{n}-1}\|^3_{H^1}
         +\varepsilon^2 \tau \big\| (\widehat{g(\psi^\varepsilon_{\tilde{n}-1})}\big)_0 {\rm e}^{{\bf i} \tau \Delta} \psi^\varepsilon_{\tilde{n}-1}\big\|_{H^1}\\
         &+\varepsilon^2 \tau \|{\rm e}^{{\bf i}\tau \Delta} h(\psi^\varepsilon_{\tilde{n}-1})\|_{H^1}+\varepsilon R_{1,\varepsilon}(w),
     \end{align*}
     for some $C_{GN}>0$.
% \begin{align*}
 % \| \psi_{n+1}\|_{H^\sigma} \le 2 R_0&+ \tau \|(\psi^n)^2(\varphi_1(-2 {\bf i} \tau \Delta)\overline{\psi_n})\|_{H^\sigma}\\
% &\le 2 R_0+C \tau \|\psi_n\|_{H^\sigma}^3\\
% &\le 2 R_0+C \tau R_0^2\|\psi_n\|_{H^\sigma}.
%+2 \tau \| (\widehat{g(u^n)}\big)_0 {\rm e}^{{\bf i}\tau \Delta } u^n\|_{H^\sigma}
%+\tau\| h(u^n)\|_{H^\sigma}.
%\end{align*}

By definition of $g$ and using the Holder's inequality, we have
\begin{align*}
|\big(\widehat{g(\psi^\varepsilon_{\tilde{n}-1})}\big)_0|&\le \displaystyle\frac{1}{2\pi}\big[\left\|\psi_{\tilde{n}-1}^\varepsilon\right\|^2+\int_{-\pi}^{\pi}\left|\psi^\varepsilon_{\tilde{n}-1}(x)\right|\left|\varphi(-2{\bf i}\tau \Delta)\overline{\psi_{\tilde{n}-1}^\varepsilon}(x)\right| \,{\rm d}x\big]\\
&\le \displaystyle\frac{1}{2\pi}\big[\left\|\psi_{\tilde{n}-1}^\varepsilon\right\|^2+\big(\int_{-\pi}^{\pi}\left| \psi_{\tilde{n}-1}^\varepsilon(x)\right|^2\,{\rm d}x\big)^{\frac{1}{2}}
\big(\int_{-\pi}^{\pi}\left|\varphi(-2{\bf i}\tau \Delta)\overline{\psi_{\tilde{n}-1}^\varepsilon}(x)\right|^2\,{\rm d}x\big)^{\frac{1}{2}}\big]\\
&=\displaystyle\frac{1}{2\pi}\left[\left\|\psi^\varepsilon_{\tilde{n}-1}\right\|^2_{H^1}+\left\|\psi^\varepsilon_{\tilde{n}-1}\right\|_{H^1}\left\|\varphi(-2{\bf i}\tau \Delta) \overline{\psi_{\tilde{n}-1}^\varepsilon}\right\|_{H^1}\right]\\
&\lesssim  \left\| \psi_{\tilde{n}-1}^\varepsilon\right\|_{H^1}^2.
\end{align*}
This then yields
$$
\big\| (\widehat{g(\psi_{\tilde{n}-1}^\varepsilon)}\big)_0 {\rm e}^{{\bf i} \tau \Delta} \psi_{\tilde{n}-1}^\varepsilon\big\|_{H1}\le |  (\widehat{g(\psi_{\tilde{n}-1}^\varepsilon)}\big)_0|
\| {\rm e}^{{\bf i} \tau \Delta} \psi^\varepsilon_{\tilde{n}-1}\|_{H^1}\lesssim  \|\psi^\varepsilon_{\tilde{n}-1} \|^3_{H^1}.
$$
It remains to bound the term $\left\| h(\psi^\varepsilon_{\tilde{n}-1})\right\|_{H^1}$. First note that for $w \in H^1$, we have
$$
\left| \varphi_1\left(2 {\bf i} \tau \ell^2\right) \hat{w}_\ell \right|\le C \left|\hat{w}_\ell \right|, \quad \ell \in \mathbb{Z}.
$$
Then, we get
\begin{align*}
\| h(\psi^\varepsilon_{\tilde{n}-1})\|_{H^1}^2 &= \displaystyle\sum_{\ell \in\mathbb{Z}}\left(1+\left|\ell\right|\right)^{2}| (\widehat{h(\psi^\varepsilon_{\tilde{n}-1})})_\ell |^2\\
&\le  \displaystyle\sum_{\ell \in\mathbb{Z}}\left(1+\left|\ell \right|\right)^{2} \big|\big( \widehat{\psi^\varepsilon}_{\tilde{n}-1,\ell} \big)^*\widehat{\psi^\varepsilon}_{\tilde{n}-1,\ell}\widehat{\psi^\varepsilon}_{\tilde{n}-1,\ell} \big|^2\\
&\hspace{2cm}+\displaystyle\sum_{\ell \in\mathbb{Z}}\left(1+\left|\ell \right|\right)^{2} \big| \varphi\left(2 {\bf i} \tau \ell ^2\right)  \big( \widehat{\psi^\varepsilon}_{\tilde{n}-1,\ell} \big)^*\widehat{\psi^\varepsilon}_{\tilde{n}-1,\ell}\widehat{\psi^\varepsilon}_{\tilde{n}-1,\ell}\big|^2\\
&\lesssim \displaystyle\sum_{\ell\in\mathbb{Z}}\left(1+\left|\ell\right|\right)^{2} \big| \big( \widehat{\psi^\varepsilon}_{\tilde{n}-1,\ell} \big)^*\widehat{\psi^\varepsilon}_{\tilde{n}-1,\ell}\widehat{\psi^\varepsilon}_{\tilde{n}-1,\ell}\big|^2\\
&\lesssim \|\psi^\varepsilon_{\tilde{n}-1}\|^3_{H^1}.
\end{align*}
Then, we end up with
\begin{align*}
    \|\psi^\varepsilon_{\tilde{n}}\|_{H^1}&\le (1+\varepsilon)R_{1,\varepsilon}(w)+1+\varepsilon^2\tau C_{GN}  \|\psi^\varepsilon_{\tilde{n}-1}\|^3_{H^1}\\
    &\le  (1+\varepsilon)R_{1,\varepsilon}(w)+1+\varepsilon^2\tau C_{GN} (R_{1,\varepsilon}(w)+1)^3.
\end{align*}
This shows that for sufficiently small $\tau$ (independently on $\varepsilon$), there exists a random variable $R_{2,\varepsilon}(\omega)\ge R_{1,\varepsilon}(\omega)+1$ such that
$$
\|\psi^\varepsilon_{\tilde{n}}\|_{H^1}\le R_{2,\varepsilon}(\omega).
$$
This shows that $\psi^\varepsilon_n=\psi_n^{\varepsilon,R_{2,\varepsilon}(\omega)}$, for $n=1,\dots,\tilde{n}$ and 
$$
\max_{1\le n \le N_\varepsilon} \|\psi^{\varepsilon, R_{2,\varepsilon}(\omega)}(t_n)-\psi^{\varepsilon, R_{2,\varepsilon}(\omega)}_n\|\ge \varepsilon^2 \tau^\delta,
$$
provided $\tau$ sufficiently small. In particular, this means that one can always find a sufficiently large $k$ such that 
$$
\max_{1\le n \le N^k_\varepsilon} \|\psi^{\varepsilon, R_{2,\varepsilon}(\omega)}(t_n)-\psi^{\varepsilon, R_{2,\varepsilon}(\omega)}_n\|\ge \varepsilon^2 \tau_k^\delta,
$$
being $\{\tau_k\}$ the subsequence defined in Lemma \ref{long_term_cub_trunc}. By \eqref{long_term_path_trunc_cub}, this is impossible and then a contradiction holds. The proof is concluded.
\end{proof}

\end{document}
